\documentclass[tikz,border=0pt]{standalone}
% Shared visual language for every block diagram in figures/diagrams.
% Build with XeLaTeX (see figures/diagrams/build.sh).
\usepackage{fontspec}
\setmainfont{Inter}[Numbers=Lining]
\setmonofont{DejaVu Sans Mono}[Scale=0.86]
\usepackage{xcolor}
\usetikzlibrary{arrows.meta,positioning,calc,fit,backgrounds,shapes.geometric,
  shapes.misc,matrix,decorations.pathreplacing,patterns,quotes}

% Palette (matches python/figstyle.py) --------------------------------------
\definecolor{surface}{HTML}{FCFCFB}
\definecolor{ink}{HTML}{0B0B0B}
\definecolor{ink2}{HTML}{52514E}
\definecolor{muted}{HTML}{898781}
\definecolor{hair}{HTML}{E1E0D9}
\definecolor{memline}{HTML}{2A78D6}  \definecolor{memfill}{HTML}{E7F0FC}
\definecolor{aluline}{HTML}{14875E}  \definecolor{alufill}{HTML}{E3F5EE}
\definecolor{ctlline}{HTML}{B07800}  \definecolor{ctlfill}{HTML}{FDF2D9}
\definecolor{regline}{HTML}{6250D6}  \definecolor{regfill}{HTML}{ECEAFB}
\definecolor{wire}{HTML}{3D3C39}
\definecolor{accent}{HTML}{EB6834}   \definecolor{accentfill}{HTML}{FDE8DF}
\definecolor{pos}{HTML}{E34948}      \definecolor{posfill}{HTML}{FBE1E0}
\definecolor{neg}{HTML}{2A78D6}      \definecolor{negfill}{HTML}{DCE9FA}
\definecolor{zerofill}{HTML}{F0EFEC}

\tikzset{
  every picture/.style={line cap=round, line join=round, font=\small, text=ink},
  % canvas
  canvas/.style={show background rectangle,
    background rectangle/.style={fill=surface, draw=hair, rounded corners=8pt, line width=0.6pt},
    inner frame xsep=16pt, inner frame ysep=14pt},
  ttl/.style={font=\bfseries\large, anchor=north west, inner sep=0pt, text=ink},
  sub/.style={font=\small, anchor=north west, inner sep=0pt, text=ink2},
  % blocks
  blk/.style={rounded corners=3pt, line width=0.8pt, align=center, minimum height=9mm,
    inner xsep=5pt, inner ysep=3pt},
  mem/.style={blk, draw=memline, fill=memfill},
  alu/.style={blk, draw=aluline, fill=alufill},
  ctl/.style={blk, draw=ctlline, fill=ctlfill},
  reg/.style={blk, draw=regline, fill=regfill, rounded corners=1pt},
  io/.style={blk, draw=muted, fill=white},
  ghost/.style={blk, draw=muted, dashed, fill=surface},
  % arithmetic primitives
  add/.style={circle, draw=aluline, fill=white, line width=0.8pt, minimum size=6.2mm,
    inner sep=0pt, font=\normalsize},
  shf/.style={rectangle, rounded corners=1.5pt, draw=aluline, fill=white, line width=0.7pt,
    font=\ttfamily\footnotesize, inner xsep=2.5pt, inner ysep=2pt, minimum height=5mm},
  gain/.style={isosceles triangle, isosceles triangle apex angle=60, draw=aluline, fill=white,
    line width=0.7pt, inner sep=1.2pt, font=\footnotesize},
  dn/.style={circle, draw=wire, fill=white, line width=0.7pt, minimum size=8mm, inner sep=0pt,
    font=\footnotesize},
  mux/.style={trapezium, trapezium angle=72, shape border rotate=270, draw=ctlline, fill=ctlfill,
    line width=0.8pt, minimum height=13mm, minimum width=5mm, inner sep=1pt, font=\scriptsize},
  net/.style={draw=memline, fill=memfill, rounded corners=1.5pt, font=\footnotesize, inner xsep=3pt,
    inner ysep=1.6pt, minimum width=5.5mm, minimum height=4.8mm},
  % wires
  sig/.style={-{Latex[length=2.2mm,width=1.6mm]}, draw=wire, line width=0.75pt},
  bus/.style={-{Latex[length=2.6mm,width=2mm]}, draw=wire, line width=1.5pt},
  ctrl/.style={-{Latex[length=1.8mm,width=1.4mm]}, draw=ctlline, line width=0.65pt, dash pattern=on 3pt off 2pt},
  nohead/.style={-, draw=wire, line width=0.75pt},
  dot/.style={circle, fill=wire, inner sep=0pt, minimum size=3.2pt},
  % labels
  lbl/.style={font=\footnotesize, text=ink2, fill=surface, inner sep=1.2pt},
  lblc/.style={font=\footnotesize, text=ink2, inner sep=1.2pt},
  note/.style={font=\footnotesize, text=ink2, align=left, inner sep=0pt},
  head/.style={font=\bfseries\small, text=ink, inner sep=0pt},
  code/.style={font=\ttfamily\footnotesize, text=ink},
  group/.style={draw=hair, fill=white, rounded corners=6pt, line width=0.8pt, inner sep=7pt},
  grouplabel/.style={font=\bfseries\footnotesize, text=muted, anchor=south west, inner sep=2pt},
}

% Layers: groups sit above the canvas background and below the main layer.
\pgfdeclarelayer{groups}
\pgfsetlayers{background,groups,main}

% A register: box with a clock wedge at the bottom edge.
\newcommand{\clk}[1]{\draw[regline, line width=0.6pt] ($(#1.south)+(-1.4mm,0)$) -- ($(#1.south)+(0,1.6mm)$) -- ($(#1.south)+(1.4mm,0)$);}
% Bus width tick: \bw{(point)}{width}
\newcommand{\bw}[2]{\draw[wire, line width=0.6pt] ($#1+(-0.9mm,-1.2mm)$) -- ($#1+(0.9mm,1.2mm)$); \node[font=\scriptsize, text=ink2, above=0.4mm, inner sep=0.5pt] at #1 {#2};}
\newcommand{\var}[1]{\textit{#1}}
\newcommand{\usub}[1]{\ifcase#1 ₀\or ₁\or ₂\or ₃\or ₄\or ₅\fi}
\AtBeginDocument{\hyphenpenalty=10000\exhyphenpenalty=10000}

\usepackage{tikz-timing}
\begin{document}
\begin{tikzpicture}[canvas]
  \node[ttl] at (0,0) {Frame protocol: ready/valid handshakes};
  \node[sub] at (0,-0.62) {2 × 2 frame, ARCH 0: four pixels with one idle cycle, four compute clocks, four coefficients with one back-pressure stall.};
  \def\sx{4.0}
  \foreach \lab/\y/\seq in {
      clk/-2.40/{30{C}},
      in\_valid/-3.25/{4H 2L 4H 20L},
      in\_ready/-4.10/{10H 18L 2H},
      in\_pixel/-4.95/{2D{p₀} 2D{p₁} 2U 2D{p₂} 2D{p₃} 20U},
      state/-5.80/{10D{LOAD} 8D{SEP: N compute clocks} 10D{DRAIN} 2D{LOAD}},
      out\_valid/-6.65/{18L 10H 2L},
      out\_ready/-7.50/{20H 2L 8H},
      out\_coeff/-8.35/{18U 2D{c₀} 4D{c₁} 2D{c₂} 2D{c₃} 2U},
      out\_index/-9.20/{18U 2D{0} 4D{1} 2D{2} 2D{3} 2U},
      out\_last/-10.05/{26L 2H 2L}} {
    \node[font=\ttfamily\footnotesize, anchor=east] at (\sx-0.25,\y+0.2) {\lab};
    \begin{scope}[shift={(\sx,\y)}]
      \timing[timing/xunit=4.6mm, timing/yunit=4.2mm, timing/slope=0.08, timing/d/text/.style={font=\scriptsize},
              draw=wire, line width=0.75pt] at (0,0) {\seq};
    \end{scope}
  }
  % cycle grid and annotations
  \foreach \c in {0,...,15} \draw[hair, line width=0.4pt] (\sx+\c*0.92,-2.05) -- (\sx+\c*0.92,-10.2);
  \tikzset{ann/.style={font=\footnotesize, text=ink2, align=center}}
  \draw[accent, line width=1pt] (\sx+2*0.92,-3.1) rectangle ++(0.92,-0.75);
  \node[ann, anchor=north] at (\sx+2.5*0.92,-10.45) {idle:\\\texttt{in\_valid} = 0};
  \draw[accent, line width=1pt] (\sx+10*0.92,-7.35) rectangle ++(0.92,-1.95);
  \node[ann, anchor=north] at (\sx+10.5*0.92,-10.45) {stall: \var{c}₁ and\\index held};
  \draw[decorate, decoration={brace, amplitude=4pt}, draw=ink2] (\sx+5*0.92,-2.0) -- node[ann, above=5pt] {compute: \var{N} = 4 clocks} (\sx+9*0.92,-2.0);
  \node[note, anchor=north west, text width=8.0cm] at (\sx+15.6*0.92,-2.45) {%
    A transfer happens on a rising edge where valid and ready are both high.\\[4pt]
    \texttt{in\_ready} is high only in LOAD; \texttt{out\_valid} only in DRAIN.\\[4pt]
    \texttt{out\_index} is the packed row-major address, so the sink knows where each coefficient belongs; \texttt{out\_last} marks index \var{N}−1.\\[4pt]
    The testbench drives the same pattern (an idle cycle every seventh pixel, a stall every fifth output cycle) on six frames per architecture.};
\end{tikzpicture}
\end{document}
